\subsection{Generalization through the full optimizer state}
\label{app:generalization}

Generalization requires a sampling model beyond smooth optimization. In this
subsection $S=(z_1,\ldots,z_n)$ is an i.i.d. training set, $S^{(i)}$ replaces
one observation independently, $L(W)=\mathbb E_z\ell(W,z)$, and
$\widehat L_S(W)=n^{-1}\sum_i\ell(W,z_i)$. These assumptions do not
directly describe transductive node classification on a single dependent
graph or nonstationary continual-learning evaluation. Applying the bounds
to such settings requires a separate statistical model.

The augmented stored state is
\begin{equation}
 X_t=(W_t,U_t,M_t,A_t,H_t),\qquad A_t=\sqrt{V_t}.
 \label{eq:augmented-state}
\end{equation}
We retain the \emph{basis} $U_t$ with the algorithm's QR gauge, rather than
identifying all bases on the quotient. Lemma~\ref{lem:moments} shows that
coordinatewise Adam does not define a gauge-invariant state transition on
the quotient alone. Distances below compare stored frames in a common
regular chart. For fixed positive weights $\omega_j$, let
\begin{equation}
 d_\omega(X,X')=\|W-W'\|+\omega_U\|U-U'\|
 +\omega_M\|M-M'\|+\omega_A\|A-A'\|+\omega_H\|H-H'\|.
\end{equation}
Multiple routed layers are handled by adding their terms. The metric must
include all persistent fallback states as well; omitting a full Adam
fallback from the sensitivity calculation would analyze a different method.

\begin{assumption}[Regular full-state update maps]
\label{ass:stability}
For every possible training example and every pair of reachable states,
the one-example update $\Psi_t$ satisfies
\begin{align}
 d_\omega(\Psi_t(X,z),\Psi_t(X',z))&\le L_t d_\omega(X,X'),\label{eq:same-sample}\\
 d_\omega(\Psi_t(X,z),\Psi_t(X,z'))&\le b_t.\label{eq:changed-sample}
\end{align}
The constants $L_t,b_t$ are deterministic and uniform over the data sets
under comparison, and initialization is data independent and coupled.
The loss is $L_\ell$-Lipschitz in $W$ on the reachable parameter region.
\end{assumption}

This is a fixed-configuration analysis. If hyperparameters are chosen by
validation, the entire selection procedure becomes part of the learning
algorithm. A neighboring observation can change its discrete choice;
fixed-configuration stability does not automatically establish stability
of that selection pipeline. A separate selection argument or uniform
control of its branches is needed for that stronger conclusion.

\begin{theorem}[Full-state stability and expected generalization]
\label{thm:generalization}
Under Assumption~\ref{ass:stability}, couple the sampling indices of the
two runs so that the replaced observation is used with marginal probability
$1/n$ at each step. Then
\begin{align}
 \mathbb E d_\omega(X_T,X_T^{(i)})
 &\le\frac1n\sum_{t<T}b_t\prod_{j=t+1}^{T-1}L_j,\label{eq:stability-product}\\
 \left|\mathbb E_{S,\Omega}\{L(W_T)-\widehat L_S(W_T)\}\right|
 &\le\beta_T:=\frac{L_\ell}{n}
                    \sum_{t<T}b_t\prod_{j=t+1}^{T-1}L_j.
 \label{eq:expected-generalization}
\end{align}
Here $\Omega$ denotes the algorithm's randomness. If the loss belongs to
$[0,M_\ell]$, with probability at least $1-\delta$ over $S$,
\begin{equation}
 \left|\mathbb E_\Omega[L(W_T)-\widehat L_S(W_T)]\right|
 \le\beta_T+(2n\beta_T+M_\ell)
                    \sqrt{\frac{\log(2/\delta)}{2n}}.
 \label{eq:stability-highprob}
\end{equation}
The high-probability statement concerns the randomness-averaged output;
it is not automatically a guarantee for one realized seed.
\end{theorem}

\begin{proof}
Use the same sampled index and all other random choices in both runs.
The triangle inequality, \eqref{eq:same-sample}, and
\eqref{eq:changed-sample} imply pathwise
\begin{equation}
 d_\omega(X_{t+1},X_{t+1}^{(i)})
 \le L_t d_\omega(X_t,X_t^{(i)})
           +b_t\mathbf1\{I_t=i\}.
\end{equation}
Taking expectation and using the marginal sampling probability gives the
recursion with additive term $b_t/n$. Uniformity of $L_t,b_t$ is what
allows this unconditional step; it does not require independence between
the past state and the indicator. Starting from zero distance and iterating
proves \eqref{eq:stability-product}. Lipschitzness in $W$ gives a loss
stability bound $\beta_T$ for every evaluation example.

For completeness, introduce an independent ghost observation $z_i'$ and
the replacement set $S^{(i)}$. Exchangeability gives
\begin{align}
 \mathbb E\{L(W_T(S))-\widehat L_S(W_T(S))\}
 =\frac1n\sum_i\mathbb E\{
 \ell(W_T(S),z_i')-\ell(W_T(S^{(i)}),z_i')\}.
\end{align}
Each summand has absolute value at most $\beta_T$, proving the expected
bound. This is the standard replacement argument of algorithmic stability
\cite{bousquet2002stability,hardt2016stability}.

Let $F(S)=\mathbb E_\Omega[L(W_T)-\widehat L_S(W_T)]$.
Replacing one example changes its population term by at most $\beta_T$.
Its empirical term changes by at most $\beta_T+M_\ell/n$: unchanged
examples contribute stability, and the replaced loss is bounded in an
interval of length $M_\ell$. Thus each bounded difference is at most
$2\beta_T+M_\ell/n$. McDiarmid's inequality applied to $F(S)$, together
with $|\mathbb EF(S)|\le\beta_T$, proves
\eqref{eq:stability-highprob}.
\end{proof}

For a mini-batch of $B$ sampled indices the probability of touching the
replaced observation is at most $B/n$. The corresponding update difference
must be computed for one changed member of the averaged batch. It cannot
be assumed to equal $b_t/B$ after nonlinear Adam normalization. The same
proof applies with that batch-specific injection constant. Likewise,
random reshuffling has a uniform marginal index but does not give the
conditional unbiased-gradient assumption of
Theorem~\ref{thm:convergence} without further analysis. The two theorems
have deliberately different sampling requirements.

\subsection{Checking the stability constants for the implemented primitives}

Assumption~\ref{ass:stability} is meaningful only if its constants are
connected to the algorithm. The following bounds provide that connection
and make the sources of potentially vacuous constants explicit.

\begin{lemma}[Lipschitz building blocks]
\label{lem:building-blocks}
Suppose single-example gradients are $L_g$-Lipschitz and bounded by
$G_0$, stored moments are bounded, $\varepsilon_c,\varepsilon>0$, and
all compared QR/polar operations stay in a fixed regular region.
The normalized covariance map obeys
\begin{equation}
 \left\|\frac{G^\top G}{\|G\|^2+\varepsilon_c}
       -\frac{G'^\top G'}{\|G'\|^2+\varepsilon_c}\right\|
 \le\frac{2}{\sqrt{\varepsilon_c}}\|G-G'\|.\label{eq:cov-lip}
\end{equation}
For orthonormal $U,U'$ and bounded $C,C'$,
\begin{equation}
 \|\Pi_U(CU)-\Pi_{U'}(C'U')\|
 \le4C_0\|U-U'\|+\|C-C'\|.\label{eq:tangent-lip}
\end{equation}
The elementwise root-EMA map is jointly Lipschitz:
\begin{equation}
 \|\sqrt{\beta A^{\odot2}+(1-\beta)Z^{\odot2}}
       -\sqrt{\beta A'^{\odot2}+(1-\beta)Z'^{\odot2}}\|
 \le\sqrt\beta\|A-A'\|+\sqrt{1-\beta}\|Z-Z'\|.
 \label{eq:root-ema}
\end{equation}
For a fixed orthogonal $Q$, diagonal covariance transport is nonexpansive
in root amplitudes:
\begin{equation}
 \|\sqrt{A^{\odot2}Q^{\odot2}}-
                 \sqrt{A'^{\odot2}Q^{\odot2}}\|\le\|A-A'\|.
 \label{eq:root-transport}
\end{equation}
Together with \eqref{eq:adam-lipschitz}, these inequalities yield finite
full-state constants on the stated region.
\end{lemma}

\begin{proof}
Differentiating $G^\top G/(\|G\|^2+\varepsilon_c)$ in direction $E$
bounds its derivative by
\begin{equation}
 \left\{\frac{2\|G\|}{\|G\|^2+\varepsilon_c}
 +\frac{2\|G\|^3}{(\|G\|^2+\varepsilon_c)^2}\right\}\|E\|
 \le\frac{2}{\sqrt{\varepsilon_c}}\|E\|.
\end{equation}
Integrating along the line segment proves \eqref{eq:cov-lip}.
For \eqref{eq:tangent-lip}, expand $CU-UU^\top CU$ into differences
of its factors and use
$\|UU^\top-U'U'^\top\|_{\rm op}\le2\|U-U'\|_{\rm op}$.
The $C$ difference has coefficient one after grouping
$(I-P)(C-C')U$; the $U$ differences have total coefficient at most
$4C_0$. To prove \eqref{eq:root-ema}, view each coordinate as the
Euclidean norm of $(\sqrt\beta A_{ij},\sqrt{1-\beta}Z_{ij})$ and
apply the reverse triangle inequality followed by Minkowski.
Similarly, the transported amplitude in coordinate $j$ of a row is
the norm of $(A_kQ_{kj})_k$. Squaring the reverse-triangle bound and
summing over $j$ uses $\sum_jQ_{kj}^2=1$ and proves
\eqref{eq:root-transport}.

QR with a fixed positive-diagonal convention is smooth on full-column-rank
matrices. The polar factor is smooth when its input singular values stay
away from zero. Bounded derivatives on the stipulated compact regular
region give their finite constants. Multiplication and projection are
Lipschitz on bounded sets; the residual scale uses Lipschitz norms,
a positive denominator offset, and nonexpansive scalar clipping. Finally,
the Adam division bound is \eqref{eq:adam-lipschitz}.
\end{proof}

One can compose these bounds as a nonnegative sensitivity matrix $K_t$.
Its entries bound the output block differences
$(\delta W,\delta U,\delta M,\delta A,\delta H)$ in terms of their
input differences. For example, the fresh routed-coordinate difference is
at most $L_g\delta W+G_0\delta U$, the first-moment EMA contributes
$\beta_1\delta M+(1-\beta_1)(L_g\delta W+G_0\delta U)$ before
alignment, and the parameter update contributes
$\delta W+\eta_t\delta D$. Projection transport adds
$2H_0\delta U$ to the transported momentum comparison. QR and polar
Lipschitz factors then propagate these differences through basis and
moment alignment. If $D_\omega=\operatorname{Diag}(1,\omega_U,
\omega_M,\omega_A,\omega_H)$, a sufficient choice is
\begin{equation}
 L_t=\|D_\omega K_tD_\omega^{-1}\|_1.
 \label{eq:sensitivity-matrix}
\end{equation}
The changed-example constant is obtained by repeating this composition
with zero initial state differences and gradient difference at most
$2G_0$. This procedure includes the moment states rather than treating
Adam as a memoryless gradient map.

Very small denominator offsets can make these constants large. Moreover,
if $L_t\le1+c(\eta_t+\rho_t)$ is established in a chosen weighted
metric, the product bound is at most
$\exp\{c\sum_t(\eta_t+\rho_t)\}$. A constant tracking step over a
long horizon can therefore yield a vacuous bound. We report the product
form within this scope; we do not infer a numerical certificate for the
experiments without measuring valid constants. Identical deterministic
QR, transport, and reset operations in coupled runs do not create an
additive error independent of the sample change. Such a term would
incorrectly predict different trajectories even when the two data sets
and initializations are identical.
This does not make a thresholded reset Lipschitz: crossing its threshold
can create a discontinuity. Such a reset requires an additional branch
analysis or a restriction that both coupled runs use the same branch.
The corrected diagonal-transport update avoids that particular switch.

\begin{proposition}[A nonvacuous contractive regime and the coupling condition]
\label{prop:small-gain}
Suppose, after including or validly eliminating contractive moment states,
the coupled same-example distances obey
\begin{equation}
 \begin{pmatrix}w_+\\u_+\end{pmatrix}
 \le\begin{pmatrix}1-\eta a&\eta b\\\rho c&1-\rho d\end{pmatrix}
       \begin{pmatrix}w\\u\end{pmatrix},\qquad
 a,d>0,\ b,c\ge0,\ ad>bc,\label{eq:two-block-stability}
\end{equation}
with nonnegative diagonal entries. For $b,c>0$ choose
\begin{equation}
 \frac{\eta b}{\rho d}<\omega<\frac{\eta a}{\rho c},\qquad
 \kappa=\min\{\eta a-\omega\rho c,\ \rho d-\eta b/\omega\}>0.
\end{equation}
If replacing the sampled example adds at most $(\eta B_W,\rho B_U)$
to these two blocks, then
\begin{equation}
 \beta_T\le\frac{L_\ell(\eta B_W+\omega\rho B_U)}{n\kappa}
                \{1-(1-\kappa)^T\}.\label{eq:contractive-stability}
\end{equation}
The cases $b=0$ or $c=0$ follow by the corresponding one-sided choice
of $\omega$ or by a limit.
\end{proposition}
\begin{proof}
Multiplying the second row of \eqref{eq:two-block-stability} by $\omega$
and adding the first gives
$w_++\omega u_+\le(1-\kappa)(w+\omega u)$.
The expected changed-example injection is
$(\eta B_W+\omega\rho B_U)/n$.
Summing the geometric recursion and applying the loss Lipschitz bound
proves \eqref{eq:contractive-stability}.
\end{proof}

This proposition identifies a sufficient nonvacuous regime: each block
must contract and their cross-coupling must satisfy $ad>bc$. A routing
eigengap can improve $d$ in a genuinely locally convex neighborhood, but
PL over the larger normal neighborhood is insufficient to assert
\eqref{eq:two-block-stability}. Strong convexity of the parameter objective
in the relevant metric is also an additional condition, seldom globally
true in deep learning. For explicit moment blocks, the same argument uses
a positive weight vector with weighted column sums of their full
sensitivity matrix strictly below one. These statements give checkable
sufficient conditions, rather than claiming all empirical runs satisfy them.

\subsection{A complementary PAC-Bayes guarantee for data-dependent routing}
\label{app:pacbayes}

Stability is algorithm dependent and can become loose over long horizons.
A second route uses a posterior over the final weights, with no
contractivity assumption. It establishes a statistical guarantee for a
randomized predictor and separates optimizer-state dimension from model
capacity. We provide the elementary exponential-moment proof so that its
statistical content and its assumptions are explicit.

\begin{proposition}[PAC-Bayes bound and Grassmann-shaped perturbations]
\label{prop:pacbayes}
Let $0\le\ell(w,z)\le1$, let $\mathcal P$ be any data-independent
prior on weights, and fix $\lambda>0$ before observing $S$. With
probability at least $1-\delta$, simultaneously for every posterior
$\mathcal Q$ absolutely continuous with respect to $\mathcal P$,
\begin{equation}
 L(\mathcal Q)\le\widehat L_S(\mathcal Q)
 +\frac{\operatorname{KL}(\mathcal Q\|\mathcal P)+\log(1/\delta)}{\lambda}
 +\frac{\lambda}{8n}.\label{eq:pac-bayes}
\end{equation}
For one matrix layer, use a data-independent center $W_0$ and
\begin{align}
 \mathcal P&=\mathcal N(\operatorname{vec}W_0,s_0^2I_{md}),\\
 \mathcal Q_S&=\mathcal N(\operatorname{vec}W_T,
                   I_m\otimes(s_\parallel^2P_T+s_\perp^2(I-P_T))),
 \label{eq:routing-posterior}
\end{align}
where vectorization is rowwise and all variances are strictly positive.
The final, data-dependent projector $P_T$ is allowed. Then
\begin{align}
 \operatorname{KL}(\mathcal Q_S\|\mathcal P)
 &=\frac{\|W_T-W_0\|^2}{2s_0^2}\notag\\
 &\quad+\frac m2\left[r\,k\!\left(\frac{s_\parallel^2}{s_0^2}\right)
 +(d-r)\,k\!\left(\frac{s_\perp^2}{s_0^2}\right)\right],
 \label{eq:gaussian-kl}\\
 k(x)&=x-1-\log x.
\end{align}
\end{proposition}

\begin{proof}
For fixed $w$, Hoeffding's exponential-moment inequality for the independent
bounded losses gives
\begin{equation}
 \mathbb E_S\exp\{\lambda(L(w)-\widehat L_S(w))\}
 \le\exp(\lambda^2/(8n)).
\end{equation}
Integrating over the prior and applying Markov's inequality shows that,
with probability $1-\delta$,
\begin{equation}
 \log\mathbb E_{w\sim\mathcal P}
       e^{\lambda(L(w)-\widehat L_S(w))}
 \le\lambda^2/(8n)+\log(1/\delta).
\end{equation}
For any posterior, Jensen applied after change of measure gives
\begin{equation}
 \lambda\mathbb E_{\mathcal Q}(L-\widehat L_S)
 \le\operatorname{KL}(\mathcal Q\|\mathcal P)
      +\log\mathbb E_{\mathcal P}e^{\lambda(L-\widehat L_S)}.
\end{equation}
Combining the last two displays proves \eqref{eq:pac-bayes}; the
high-probability event is independent of the posterior, giving simultaneous
validity. This is a standard PAC-Bayesian change-of-measure argument
\cite{mcallester1999pac,seeger2002pacbayes}.

For the Gaussian specialization, the covariance in each row has eigenvalue
$s_\parallel^2$ with multiplicity $r$ and $s_\perp^2$ with multiplicity
$d-r$. Insert its trace and determinant in the Gaussian KL identity
\begin{equation}
 \operatorname{KL}=\tfrac12\left[
 \frac{\|W_T-W_0\|^2+\operatorname{tr}\Sigma_Q}{s_0^2}
 -md+md\log s_0^2-\log\det\Sigma_Q\right].
\end{equation}
Collecting the two eigenvalue groups proves \eqref{eq:gaussian-kl}.
For independent layerwise Gaussian factors, KL adds across layers.
\end{proof}

The displacement term can be related to the optimizer path without
assuming the final parameter is low rank. Orthogonality of the two
implemented branches gives
\begin{align}
 \|D_t\|^2&=\|D_t^\parallel\|^2+
                  \gamma_t^2\|G_t(I-P_t)\|^2,\\
 \|W_T-W_0\|&\le\sum_{t<T}\eta_t\|D_t\|.
 \label{eq:displacement-path}
\end{align}
For adaptive branch-to-residual scaling, $\gamma_t$ here includes that
actual scalar. Equations~\eqref{eq:tail-exact} and
\eqref{eq:displacement-path} connect raw rank tails and routing quality
to an upper bound on the path's squared length, under the requisite
conditional measurability. They do not force a smaller path or smaller
generalization gap than AdamW.

Several qualifications prevent misleading rank-based interpretations.
First, $k(x)\ge0$ and diverges as $x\downarrow0$; concentrating posterior
variance in only a few directions can \emph{increase} the KL cost in the
$m(d-r)$ residual coordinates. Second, reducing those variances may improve
the empirical perturbation loss, producing a trade-off which must be
measured. Third, the prior cannot be centered at a data-dependent fitted
weight without an additional valid prior-selection argument. The isotropic
prior above allows a data-dependent posterior orientation precisely because
\eqref{eq:pac-bayes} holds for all posteriors; no assumption that $P_T$ is
independent of the training set has been made.

The empirical term is the loss of the Gaussian-perturbed predictor, not
the unperturbed network's training loss. A certificate requires evaluating
that expectation, with a separate Monte Carlo error allowance if sampled.
Unbounded cross-entropy must be replaced by a specified bounded surrogate
or analyzed with an appropriate unbounded-loss theorem. Also, optimizing
$\lambda$ after observing the posterior is not justified by a bound proved
for one fixed value. A finite grid of $J$ predeclared values is valid after
replacing $\log(1/\delta)$ by $\log(J/\delta)$ and taking a union bound.
Neither the stability bound nor PAC-Bayes is added to the other as a
spurious improvement; they are alternative guarantees with distinct
assumptions and output objects.

\subsection{Implications that can be tested, and claims that do not follow}

The analysis suggests four concrete measurements. First, estimate the full
covariance spectrum and the distance to an \emph{independent} oracle
subspace; $GU$ alone does not reveal $\lambda_{r+1}$, and consecutive
bases do not reveal the oracle radius. Second, separate normalized and raw
covariances when comparing tracking and residual energy. Third, log actual
parameter-direction alignment or its comparison-metric defect before
claiming the rate in Theorem~\ref{thm:convergence}. Fourth, estimate
neighboring-data sensitivity with shared initialization and sampling,
including moment states, before interpreting a train--test gap as a
stability measurement. Small single-seed changes in accuracy cannot
establish any of these theoretical hypotheses.

The four main result families provide a conditional chain from spectral separation
and angles to routing quality, from a controlled tracking process to an
explicit perturbation budget, from that budget to stationarity, and from
full-state sensitivity or posterior complexity to generalization. They
do not imply universal optimizer superiority, a compulsory rank-tail error
floor with a nonzero residual branch, full gauge invariance of Adam, or a
nonvacuous numerical generalization certificate for the present deep
networks. These boundaries are part of the mathematical conclusions.
Theorem~\ref{thm:gap-collapse} additionally shows why the separated tracking
and stationarity regimes need not hold simultaneously throughout training;
Theorem~\ref{thm:tracking-PAC} supplies an explicit finite-window statistical
connection with a valid independent prior.
